% ============================================================
%  \subsection{UI/UX Design}
%  aBridgeAI — main-body summary; full screen-by-screen detail
%  is in Appendix \ref{appendix:uiux}.
% ============================================================

\subsection{UI/UX Design}

% ────────────────────────────────────────────────────────────
\subsubsection{Design Philosophy}
% ────────────────────────────────────────────────────────────

aBridgeAI adopts a \textbf{light-first, academic} design language inspired by enterprise dashboards (Metronic Demo 6) --- dense rows, crisp 1\,px borders, restrained 4--10\,px corner radii. The interface deliberately reads as a serious learning tool rather than a consumer product. Four guiding principles shape every screen:

\begin{enumerate}
  \item \textbf{Role-based clarity} --- students, instructors, and administrators see distinct navigation structures and role-appropriate destinations; the same shell component renders all three contexts from a per-role \path|navItems| prop.
  \item \textbf{Progressive disclosure} --- AI features (quiz generation, processing queue, cost telemetry) are surfaced contextually inside their host screens, only when the user has the required role and prerequisites.
  \item \textbf{Accessibility-first typography} --- body copy uses \textit{Be Vietnam Pro}, a humanist sans-serif with full Vietnamese diacritic support, paired with a 14\,px base size and 1.5 line-height for sustained reading.
  \item \textbf{Responsive shell architecture} --- a persistent \path|AppShell| renders a desktop left rail that can collapse from 256\,px (\path|w-64|) to 64\,px (\path|w-16|) by user toggle. On narrow screens the rail is replaced by an off-canvas navigation menu opened from the mobile top bar.
\end{enumerate}

The implementation uses \textbf{shadcn/ui} components on top of \textbf{Tailwind CSS v4}; the colour system is defined as CSS custom properties (\path|var(--primary)|, \path|var(--surface)|, \dots) so that the light and dark themes share a single component layer with full parity.

\paragraph{Internationalisation (i18n).} Every authenticated page exposes a globe icon in the \path|ContentTopBar| that opens a language selector; the active language is shown next to the icon (\textit{e.g.}\ ``EN''). Locale strings live in \path|src/i18n/locales/<lang>.json| --- English and Vietnamese --- with English as the fallback, and a test asserts that the two catalogues carry the same key set, so a translation added on one side but not the other fails the build rather than degrading silently at runtime. Navigation labels and form fields reference an \path|i18nKey| (such as \path|nav.dashboard|, \path|nav.career_paths|) defined in \path|src/lib/navigation.ts|, so adding a new locale requires only a JSON dictionary --- no component changes. The platform was developed with English and Vietnamese as first-class locales, matching the choice of Be Vietnam Pro as the primary typeface.

% ────────────────────────────────────────────────────────────
\subsubsection{Design Tokens}
% ────────────────────────────────────────────────────────────

Table~\ref{tab:design_tokens_main} consolidates the design tokens that govern the visual system. Full per-rule rationale and the complete spacing, layering, and interaction rules are documented in Appendix~\ref{appendix:uiux}.

\begin{longtable}{|p{0.22\textwidth}|p{0.30\textwidth}|p{0.36\textwidth}|}
  \caption{Core design tokens.} \label{tab:design_tokens_main} \\
  \hline
  \textbf{Token} & \textbf{Value} & \textbf{Usage} \\
  \hline
  \endfirsthead
  \multicolumn{3}{c}{\small\tablename~\thetable{} \textit{(continued)}} \\
  \hline
  \textbf{Token} & \textbf{Value} & \textbf{Usage} \\
  \hline
  \endhead
  \hline
  \endfoot
  Page background & \path|#f8fafc| (slate-50) & Page surface in light mode. \\
  \hline
  Surface / Card & \path|#ffffff| & Cards, panels, dialogs. \\
  \hline
  Primary accent & \path|#1e40af| (blue-800) & Primary actions, active nav, links, chart series 1. \\
  \hline
  Secondary accent & \path|#3b82f6| (blue-500) & Secondary buttons, focus rings, chart series 2. \\
  \hline
  CTA / accent & \path|#f59e0b| (amber-500) & High-attention CTAs and AI affordances. \\
  \hline
  Success & \path|#16a34a| (green-600) & Published / completed states. \\
  \hline
  Warning & \path|#ea580c| (orange-600) & Drafts, pending review. \\
  \hline
  Danger & \path|#dc2626| (red-600) & Failed jobs, destructive actions. \\
  \hline
  Border & \path|#e2e8f0| (slate-200) & 1\,px borders on cards, inputs, tables. \\
  \hline
  Typography & Be Vietnam Pro (system-ui fallback) & All body and UI text; full Vietnamese diacritic support. \\
  \hline
  Component library & shadcn/ui + Radix UI + Lucide React & Dialogs, dropdowns, sheets, tooltips, icons. \\
  \hline
  Charts & Recharts & Bar, line, area; reuses palette for chart series 1--5. \\
  \hline
  i18n keys & \path|src/lib/navigation.ts| + \path|src/i18n/locales/<lang>.json| & Translatable nav labels and copy. \\
  \hline
  Sidebar widths & 64\,px collapsed / 256\,px expanded & \path|w-16| / \path|w-64| on desktop; narrow screens use the off-canvas mobile navigation menu. \\
  \hline
  Border radius (card) & 8\,px (\path|rounded-lg|) & 10\,px hard cap project-wide. \\
  \hline
  Base spacing unit & 4\,px & All spacing is a multiple of 4\,px. \\
  \hline
  Max modal depth & 1 stacked modal & Prevents nested-modal anti-patterns. \\
  \hline
\end{longtable}

% ────────────────────────────────────────────────────────────
\subsubsection{Key Screens}
% ────────────────────────────────────────────────────────────

This subsection sketches the principal screens by role. Per-screen layouts and interaction details are illustrated in Appendix~\ref{appendix:uiux}.

\paragraph{Public surfaces.} The \textbf{Landing Page} (\path|/|) is the marketing entry point; it carries a hero, stats strip, category grid, and featured-courses carousel under a dedicated \path|TopNavBar| that is unique to this route. The \textbf{Login} screen (\path|/login|) offers Google SSO as the single sign-in path; on success, accounts with MFA enabled are routed to \path|/login/mfa| for a TOTP or recovery-code challenge before a session is issued. A \path|next| query parameter preserves deep links across the auth round-trip.

\paragraph{Student.} The \textbf{Student Dashboard} (\path|/dashboard|) keeps the layout calm: stat cards summarise enrolled courses, in-progress lessons, and recently earned milestones; a recommendations strip suggests the next course drawn from the career paths the student currently holds. The \textbf{Course Catalog} (\path|/courses|) opens with a sticky search and filter bar (query, category, difficulty); enrolled courses display an inline progress ring so students can pick up where they left off. The \textbf{Course View} (\path|/courses/$slug/learn|) splits into a sticky curriculum accordion on the left and the active lesson's content on the right (video, reading, attached resources, embedded quiz launcher). The \textbf{Quiz Session} (\path|/courses/$slug/quiz/$quizId|) is a focused, distraction-free interface with a countdown timer and a question-navigation grid; on submission students land on the attempt-review screen with correct answers and per-question explanations.

\paragraph{Instructor.} The \textbf{Teacher Dashboard} (\path|/teacher|) summarises authoring activity (total / published / draft / AI-enabled) and surfaces the most recently edited course with a one-click jump into management. The \textbf{Course Manage} screen (\path|/teacher/courses/$courseId|) is the authoring core: modules render as drag-handle cards, expanding to reveal ordered lesson and quiz items, each with a publish-state badge and inline \textbf{+ Video / + Reading / + Quiz} quick-add buttons. The \textbf{Lesson Editor} (\path|/teacher/courses/$courseId/lessons/$lessonId|) splits into a Markdown content area with attached resources on the left and a settings sidebar (Visibility, Lesson Type, Duration, Difficulty, Prerequisites, AI Material Hub) on the right.

The \textbf{Quiz Editor} (\path|/teacher/courses/$courseId/quizzes/$quizId|) is structured around three sub-tabs --- \textit{Questions}, \textit{Settings}, and \textit{Preview}. The Questions tab lists every item with its type badge, per-question time limit, and approval status; bulk-action controls let the instructor apply uniform timing or bulk-approve. A persistent banner reports how many questions are pending review, enforcing a deliberate approval gate between AI generation and student exposure. An \textbf{Open generator} button launches a modal that drives the AI quiz pipeline: source-lesson selection, question count, difficulty band, question types (multiple choice, true/false, short answer, fill in the blank), generation mode (\textit{Topic} balanced spread vs.\ \textit{Coverage} guaranteeing per-section coverage), and an \textit{Append} or \textit{Replace} policy for existing questions. Generated items land in \textit{Pending Review} state and are visible to students only after explicit instructor approval.

\paragraph{Assessment surfaces.} Two screens exist to hold a learner's full attention and are deliberately stripped of the surrounding shell. The \textbf{Quiz Session} runs a countdown against the server's authoritative clock, so closing the tab neither pauses nor forfeits the attempt; where the quiz requires a camera the workspace blocks while no live video track is present, after a short grace period that absorbs a momentary device glitch. Because an attempt admits one client at a time, opening it a second time presents a conflict screen offering an explicit takeover rather than silently splitting the session in two. The \textbf{Interview Session} begins through the course interview workspace. It requires staged identity, audio, language, preparation, and readiness checks before assessed questioning and timer start. A native LiveKit room is dispatched after readiness in the target deployment; if microphone permission or voice availability is absent, the same hybrid workspace retains typed input. Results present the learner-facing verdict and qualitative remediation after asynchronous evaluation rather than exposing assessment internals.

\paragraph{Progression surfaces.} The \textbf{Study} screens (\path|/me/study|, with \path|cards-due| and \path|review|) drive the spaced-repetition loop, presenting what is due rather than asking the learner to choose. \textbf{Progress} (\path|/me/progress|), \textbf{Career Paths} (\path|/me/career-paths|) and \textbf{Learning Programmes} (\path|/me/learning-programs|) expose the learner's standing at successively wider scope, the last of these carrying pathway selection and the change-request form whose outcome a Faculty Dean decides.

\paragraph{Manager and Dean.} The management surface (\path|/management/...|) is where institutional structure is authored: learning programmes and their versions, career paths, organisational units, course enrolments, and user administration, all scoped to the operator's own organisation. The programme detail screen carries the dean's review queue, where pathway change and drop requests are acknowledged and then approved or rejected against a fixed reason vocabulary.

\paragraph{Administrator.} The \textbf{Operations} page (\path|/admin/operations|) gives platform operators visibility into the AI background-job pipeline --- material ingestion, embedding, quiz generation, validation --- with stat cards reporting the backlog, status-pill filters narrowing the table, and a per-row \textbf{Retry} on failed jobs. The \textbf{AI Costs} page (\path|/admin/ai-costs|) tracks language-model and embedding spend across a selectable range, broken down by role and by pipeline stage so that cost can be correlated with backlog health without a context switch. Beyond these sit \textbf{Organisations}, \textbf{Users}, \textbf{Audit Logs}, \textbf{Policies} and \textbf{System Configuration} (\path|/admin/settings|), the last presenting each runtime setting with its effective value, the level that supplied it, and a dry-run preview before any change is applied.

% ────────────────────────────────────────────────────────────
\subsubsection{Navigation Architecture}
% ────────────────────────────────────────────────────────────

The authenticated experience is built around a single role-agnostic \path|AppShell| component that orchestrates a fixed-position \path|SideNavBar| and a sticky \path|ContentTopBar| around the main content. The shell receives a \path|navItems| array as a prop (defined in \path|src/lib/navigation.ts|), so the same layout code serves the student, teacher, and admin route groups; role differentiation is communicated through the \path|navItems| configuration and shared design tokens. On desktop the sidebar can collapse to an icon rail; on narrow screens it is replaced by a mobile top-bar menu, keeping the main column readable without permanently consuming horizontal space.

The \path|TopNavBar| and \path|Footer| components are reserved exclusively for the landing page, ensuring a clean visual handoff between the marketing surface and the authenticated product.
